\section{Unsmoothed prime counts at arbitrary multiplicative scales}
\label{sec:pf}

The supplied article proves positive definite covariance for smoothed prime
errors at arbitrary finitely many multiplicative scales.  Its open questions
ask whether this persists for unsmoothed prime counts after the prime-square
bias is removed \cite[Section~\textup{``Zeros, Euler products and prime races''},
open question~10]{ProveItArticle}.  We prove the multiscale part of that
extension for the ordinary counting function $\pi(x;q,a)$.  The conclusion
survives removal of every fixed finite set of zero frequencies, and yields
positive and negative excursions of every centered linear contrast in
\emph{every sufficiently long logarithmic interval}.

The only hypothesis about zero locations is Dirichlet GRH, used in the precise
sense supplied by the vendored development \cite{ProveItVendor}.  Simplicity,
disjointness of different zero sets, and rational independence of ordinates
are not assumed.  Stepanov almost periodicity of normalized prime-counting
remainders is classical: Kaczorowski and Ramar\'e prove an $L^2$ version in a
more general setting \cite[Theorem~1(e)]{KaczorowskiRamare2003}.  The extension
here combines an explicit local $L^4$ argument, a transfer to actual prime
counts, and the arithmetic separation of arbitrary scales.  We do not claim
that every ordering of the uncentered prime counts occurs, or settle the
separate question about all jet derivatives.

\subsection{Statement and normalizations}
\label{subsec:pf-statement}

Fix $q\geq1$, write $G=(\mathbb Z/q\mathbb Z)^\times$ and
$n=\varphi(q)$, and use the trivial group and $n=1$ when $q=1$.
For a character $\chi$ modulo $q$, let $\chi^*$ be its primitive inducer and
let $m_\chi(\gamma)$ denote the multiplicity of the zero
$\tfrac12+i\gamma$ of $L(s,\chi^*)$.  A missing zero has multiplicity zero.
Let $\Gamma_q$ be the set of distinct ordinates occurring in this finite
family, and define
\begin{equation}
\label{eq:pf-v}
 v_a(\gamma)=\frac1n\sum_{\chi\bmod q}
       \overline{\chi(a)}m_\chi(\gamma),
 \qquad
 r_q(a)=\#\{b\in G:b^2=a\}.
\end{equation}
Thus coincidences between zeros of different functions are grouped before
any square is taken.  Conjugation gives
$v_a(-\gamma)=\overline{v_a(\gamma)}$, and $v_a(0)$ is real.
Put $\operatorname{li}(x)=\int_2^x du/\log u$, and set
\begin{equation}
\label{eq:pf-X}
 X_a(t)=t e^{-t/2}
   \left(\pi(e^t;q,a)-\frac{\operatorname{li}(e^t)}n\right)
       +\frac{r_q(a)}n+2v_a(0),\qquad t\geq\log2.
\end{equation}
The term $r_q(a)/n$ removes prime squares, while $2v_a(0)$ removes the
possible central zero contribution.  The latter is not assumed to vanish.
These constants differ from those of the article's smoothed ramp statistic
because the Mellin denominator is now $s$, rather than $s(s-1)$.

Fix distinct real shifts $h_1,\ldots,h_r$ and put $D=rn$.
They correspond to the multiplicative scales $e^{h_1},\ldots,e^{h_r}$.
For $\gamma\ne0$ define a vector in $\mathbb C^D$ by
\begin{equation}
\label{eq:pf-z}
 z(\gamma)_{j,a}
       =-\frac{e^{i\gamma h_j}v_a(\gamma)}{\tfrac12+i\gamma}.
\end{equation}
If $J\subset\Gamma_q\setminus\{0\}$ is a finite symmetric set, remove its
frequencies from the observed vector:
\begin{equation}
\label{eq:pf-deletion}
 \mathcal X_J(t)_{j,a}
       =X_a(t+h_j)-\sum_{\gamma\in J}
                         z(\gamma)_{j,a}e^{i\gamma t}.
\end{equation}
The vector is real.  Its domain starts at any fixed $u_0$ for which
$u_0+h_j\geq\log2$ for every $j$.

\begin{theorem}[Unsmoothed multiscale covariance and occupation]
\label{thm:pf-main}
Assume Dirichlet GRH for the characters modulo $q$.  For the data just
defined, the series
\begin{equation}
\label{eq:pf-cov}
 C_J=\sum_{\gamma\in\Gamma_q\setminus(J\cup\{0\})}
                   z(\gamma)z(\gamma)^*
\end{equation}
converges to a real symmetric positive definite matrix.  Uniformly over
$u\geq u_0$, as $L\longrightarrow\infty$,
\begin{equation}
\label{eq:pf-moments}
 \frac1L\int_u^{u+L}\mathcal X_J(t)\,dt\longrightarrow0,
 \qquad
 \frac1L\int_u^{u+L}\mathcal X_J(t)\mathcal X_J(t)^{\mathsf T}\,dt
                       \longrightarrow C_J.
\end{equation}
There are constants $\varepsilon,\eta,L_0>0$ such that, for every real unit
vector $w\in\mathbb R^D$, every $u\geq u_0$ and every $L\geq L_0$,
\begin{align}
\label{eq:pf-positive}
 \bigl|\{t\in[u,u+L]:w\cdot\mathcal X_J(t)\geq\varepsilon\}\bigr|
       &\geq\eta L,\\
\label{eq:pf-negative}
 \bigl|\{t\in[u,u+L]:w\cdot\mathcal X_J(t)\leq-\varepsilon\}\bigr|
       &\geq\eta L.
\end{align}
Here $|\cdot|$ denotes Lebesgue measure.  Consequently the closed convex
hull of the essential values on each such interval contains the closed
Euclidean ball of radius $\varepsilon$ about zero.

Moreover, there are $c,K,T_0>0$, depending only on $q$ and the shifts,
such that
\begin{equation}
\label{eq:pf-band}
 \lambda_{\min}(C_J)\geq\frac{cH}{T^2\log(2T)}
\end{equation}
whenever $T\geq T_0$, $K\log(2T)\leq H\leq T$, and
$T>\max_{\gamma\in J}|\gamma|$; the last maximum is zero if $J$ is empty.
\end{theorem}

The constants may depend on the fixed shifts and deteriorate when shifts
coalesce.  No uniformity as the number of coordinates grows is asserted.
In terms of $x=e^t$, the intervals above are $[e^u,e^{u+L}]$ and their
measure is $dx/x$.  Thus the theorem supplies bounded waiting on the
logarithmic scale for every centered directional contrast, with a positive
proportion of logarithmic measure on both sides.

\subsection{Local fourth moments of zero sums}
\label{subsec:pf-local}

For a locally $p$-integrable function on $\mathbb R$, write
\[
 \|f\|_{S^p}=\sup_{u\in\mathbb R}
           \left(\int_u^{u+1}|f(t)|^p\,dt\right)^{1/p}.
\]
For vector-valued functions, use the Euclidean norm inside the integral.
Convergence in this norm is uniform local $L^p$ convergence; in particular
it implies convergence as distributions.  We first prove the analytic
estimate that permits unsmoothing.

\begin{lemma}[Binned local moments]
\label{lem:pf-bins}
Let $P(t)=\sum_\lambda c_\lambda e^{i\lambda t}$ be a finite exponential
sum with real frequencies, and put
$b_j=\sum_{\lambda\in[j,j+1)}|c_\lambda|$.  There is an absolute constant
$A$ such that
\begin{equation}
\label{eq:pf-bins}
 \|P\|_{S^2}\leq A\|b\|_{\ell^2},
 \qquad
 \|P\|_{S^4}\leq A\|b\|_{\ell^{4/3}}.
\end{equation}
Consequently, if a countable locally finite frequency family has
$b\in\ell^{4/3}(\mathbb Z)$, its finite symmetric truncations converge in
$S^4$.
\end{lemma}

\begin{proof}
Consider the nonnegative function
\[
 W(s)=\left(\frac{\sin(s/2)}{s/2}\right)^2,
 \qquad W(0)=1.
\]
It has a positive lower bound on $[-1/2,1/2]$ and satisfies the exact
identity
\[
 W(s)=\int_{-1}^1(1-|\xi|)e^{i\xi s}\,d\xi.
\]
Its Fourier transform is therefore bounded and supported in $[-1,1]$,
with the usual factor $2\pi$ in the inverse convention.  Multiplying
$|P|^2$ by $W(t-u-1/2)$ and integrating over $\mathbb R$ gives
\begin{align*}
 \int_u^{u+1}|P(t)|^2\,dt
 &\leq C\int_{\mathbb R}|P(t)|^2W(t-u-1/2)\,dt\\
 &\leq C'\sum_{|\lambda-\mu|\leq1}
                             |c_\lambda||c_\mu|\\
 &\leq C''\sum_j b_j^2.
\end{align*}
Only equal or adjacent bins can interact, and
$2b_jb_{j+1}\leq b_j^2+b_{j+1}^2$.  The constants do not depend on $u$.

Apply this $S^2$ estimate to $P^2$.  Its absolute coefficient mass $d_k$
in bin $[k,k+1)$ is at most
\[
 d_k\leq(b*b)_k+(b*b)_{k-1}.
\]
Indeed, adding frequencies in bins $j$ and $\ell$ gives a frequency in
bin $j+\ell$ or $j+\ell+1$; grouping coincident sums can only decrease
absolute mass.  Young's convolution inequality yields
\[
 \|P\|_{S^4}^4
   =\|P^2\|_{S^2}^2
   \leq C\|d\|_{\ell^2}^2
   \leq4C\|b*b\|_{\ell^2}^2
   \leq4C\|b\|_{\ell^{4/3}}^4.
\]
This proves the second estimate after enlarging $A$.  Applying it to
differences of truncations proves the final assertion, since the
$\ell^{4/3}$ norm of their omitted bin masses tends to zero.
\end{proof}

For the fixed Dirichlet family, the standard local zero count gives
\begin{equation}
\label{eq:pf-localcount}
 \sum_{\gamma\in\Gamma_q\cap[j,j+1)}\sum_{\chi\bmod q}m_\chi(\gamma)
                 \ll_q\log(2+|j|).
\end{equation}
This count includes multiplicities and holds for both signs.  Therefore
the absolute coefficient masses of
\begin{equation}
\label{eq:pf-F}
 F_a(t)=-\sum_{\substack{\gamma\in\Gamma_q\\\gamma\ne0}}
                  \frac{v_a(\gamma)}{\tfrac12+i\gamma}e^{i\gamma t}
\end{equation}
obey
\begin{equation}
\label{eq:pf-bdecay}
 b_j\ll_q\frac{\log(2+|j|)}{1+|j|}.
\end{equation}
The finitely many bins near zero cause no difficulty because
$|\tfrac12+i\gamma|\geq1/2$.
Lemma~\ref{lem:pf-bins} proves convergence of \eqref{eq:pf-F} in $S^4$,
and in particular a uniform bound for its local fourth moments.  Comparing
the tails of $\sum b_j^{4/3}$ and $\sum b_j^2$ with integrals also gives,
for $Y\geq2$,
\begin{equation}
\label{eq:pf-truncation}
 \left\|F_a+
   \sum_{0<|\gamma|\leq Y}
       \frac{v_a(\gamma)e^{i\gamma t}}{\tfrac12+i\gamma}\right\|_{S^4}
       \ll_q Y^{-1/4}\log(2Y),
\end{equation}
and the corresponding $S^2$ bound is
$O_q(Y^{-1/2}\log(2Y))$.  No lower bound on gaps between ordinates has
entered these estimates.

\subsection{From the explicit formula to ordinary prime counts}
\label{subsec:pf-transfer}

We use the classical explicit formula and local zero count for primitive
Dirichlet $L$-functions, in the distributional normalization also used in
\cite{KaczorowskiRamare2003}; the fixed-argument Landau formula needed
below is recorded separately.  The standard GRH consequence
\begin{equation}
\label{eq:pf-pointwise}
 \theta(x;q,a)=\frac{x}{n}+O_q\bigl(\sqrt{x}\log^2(2x)\bigr)
\end{equation}
will be used only for deterministic remainder estimates.  These are
classical analytic inputs, not new conclusions about zero locations.

Let
\[
 \psi(x;q,a)=\sum_{\substack{p^k\leq x\\p^k\equiv a\bmod q}}\log p,
 \qquad
 \theta(x;q,a)=\sum_{\substack{p\leq x\\p\equiv a\bmod q}}\log p.
\]

\begin{lemma}[The unsmoothed expansion]
\label{lem:pf-transfer}
For each $a\in G$, there is a remainder $R_a$ such that, almost everywhere
on $t\geq\log2$,
\begin{equation}
\label{eq:pf-XF}
 X_a(t)=F_a(t)+R_a(t),
 \qquad
 \|R_a\|_{L^4([U,U+1])}\ll_q U^{-1}\quad(U\geq2).
\end{equation}
The remainder is locally bounded, belongs to $L^2([\log2,\infty))$, and
has uniformly bounded local fourth moments on that half-line.
\end{lemma}

\begin{proof}
For a primitive character $\chi$, put $\delta_\chi=1$ in the conductor-one
case and $0$ otherwise.  The explicit formula, interpreted as an identity
of distributions on $t>0$, states
\begin{equation}
\label{eq:pf-psi-character}
 e^{-t/2}\bigl(\psi(e^t,\chi)-\delta_\chi e^t\bigr)
 =-\sum_\gamma\frac{m_\chi(\gamma)e^{i\gamma t}}
                         {\tfrac12+i\gamma}
       +O_\chi\bigl((1+t)e^{-t/2}\bigr),\qquad t\geq\log2.
\end{equation}
The term denoted by $O_\chi$ consists of the constant, the contribution
at zero, and the trivial zeros.  It is smooth away from $t=0$ and has
the stated bound.  One may obtain the distributional formula by
integrating $-L'/L(s,\chi)$ against a smooth compactly supported Mellin
test and moving the contour; rapid decay of the test transform makes
the zero sum convergent.  Dirichlet GRH turns each nontrivial zero term
into the pure phase displayed in \eqref{eq:pf-psi-character}.

The bin estimates prove $S^4$, hence $L^1_{\mathrm{loc}}$, convergence
of this same zero sum.  Its limit represents the distribution in
\eqref{eq:pf-psi-character}, so the equality holds almost everywhere
for the usual arithmetic step function.  The convention at its
countably many jumps does not affect any conclusion here.  If $\chi$
is imprimitive, the finitely many removed prime-power sequences have
total size $O_q(1+t)$ before normalization.  Thus the same estimate
holds with the zeros of $\chi^*$.

Character orthogonality and the central term
$-(\tfrac12)^{-1}v_a(0)=-2v_a(0)$ now give
\begin{equation}
\label{eq:pf-psi}
 E_{\psi,a}(t):=e^{-t/2}\left(\psi(e^t;q,a)-\frac{e^t}{n}\right)
       =F_a(t)-2v_a(0)+O_q\bigl((1+t)e^{-t/2}\bigr).
\end{equation}
In particular $E_{\psi,a}$ has bounded local fourth moments uniformly
on $t\geq\log2$.

The contribution of prime squares is
\[
 \psi(x;q,a)-\theta(x;q,a)
       =\sum_{b^2=a}\theta(\sqrt{x};q,b)
                  +O_q\bigl(x^{1/3}\log^2(2x)\bigr).
\]
The error follows by bounding the contributions of prime powers of
exponent at least three.  Applying \eqref{eq:pf-pointwise} at
$\sqrt{x}$, and then dividing by $\sqrt{x}$, yields
\begin{equation}
\label{eq:pf-theta}
 E_{\theta,a}(t):=e^{-t/2}\left(\theta(e^t;q,a)-\frac{e^t}{n}\right)
 =E_{\psi,a}(t)-\frac{r_q(a)}n
                    +O_q\bigl((1+t)^2e^{-t/6}\bigr).
\end{equation}
This also has uniformly bounded local fourth moments.  Notice that
the correction is $r_q(a)/n$, not $(r_q(a)-1)/n$: we are centering each
coordinate against $\operatorname{li}(x)/n$, including the principal
direction.

Partial summation gives the exact identity
\begin{align}
\label{eq:pf-volterra}
 t e^{-t/2}\left(\pi(e^t;q,a)-\frac{\operatorname{li}(e^t)}n\right)
 &=E_{\theta,a}(t)+V_a(t)+\frac{2t e^{-t/2}}{n\log2},\\
 V_a(t)&=t e^{-t/2}\int_{\log2}^t
                    \frac{e^{v/2}E_{\theta,a}(v)}{v^2}\,dv.
 \nonumber
\end{align}
We claim that $\|V_a\|_{L^4([U,U+1])}\ll_q U^{-1}$.
Split the integral at $v=t/2$.  On the earlier part,
\eqref{eq:pf-pointwise} gives $E_{\theta,a}(v)\ll_q(1+v)^2$, and the
resulting contribution is $O_q(t^2e^{-t/4})$.  On the later part put
$s=t-v$ and use $t/(t-s)^2\leq4/t$.  For $t\in[U,U+1]$,
\[
 |V_{a,\mathrm{late}}(t)|
 \leq\frac4U\int_0^{(U+1)/2}e^{-s/2}|E_{\theta,a}(t-s)|\,ds.
\]
For large $U$, all translated unit intervals on the right lie in the
half-line where the local fourth moments are bounded.  Minkowski's
inequality, followed by $\int_0^\infty e^{-s/2}ds=2$, proves the claim.

Combining \eqref{eq:pf-psi}, \eqref{eq:pf-theta}, and
\eqref{eq:pf-volterra} proves \eqref{eq:pf-XF} with exactly the
centering in \eqref{eq:pf-X}.  The remainders constructed in these
steps are locally bounded.  Finally
\[
 \int_U^{U+1}|R_a(t)|^2\,dt
       \leq\left(\int_U^{U+1}|R_a(t)|^4\,dt\right)^{1/2}
       \ll_q U^{-2},
\]
and summation over integer $U$ proves global square integrability.
\end{proof}

\subsection{Uniform covariance and arithmetic separation of scales}
\label{subsec:pf-covariance}

\begin{lemma}[Uniform long-window moments]
\label{lem:pf-moments}
Let $F:\mathbb R\to\mathbb R^D$ be the $S^2$ limit of finite
exponential sums with no constant frequency, and write its grouped
coefficients as $z_\gamma$, with
$z_{-\gamma}=\overline{z_\gamma}$ and $\sum_\gamma\|z_\gamma\|^2<\infty$.
Then its long-window mean is zero and its long-window second moment is
$\sum_\gamma z_\gamma z_\gamma^*$, uniformly over the window origin.
The same limits hold on a half-line for $F+R$ if $R$ belongs to $L^2$
on that half-line.
\end{lemma}

\begin{proof}
For a finite sum, the assertion follows by expanding products and using
\[
 \left|\frac1L\int_u^{u+L}e^{i\omega t}\,dt\right|
                 \leq\frac2{L|\omega|}\qquad(\omega\ne0),
\]
independently of $u$.  If $P$ approximates $F$ in $S^2$, covering an
interval of length $L\geq1$ by at most $L+1$ unit intervals gives
\[
 \frac1L\int_u^{u+L}\|F-P\|^2\,dt\leq2\|F-P\|_{S^2}^2.
\]
Cauchy--Schwarz controls the errors in the first and second moments;
the local second-moment bounds of $F$ and $P$ stay bounded during the
approximation.  Pass first to large $L$ and then to the $S^2$ limit.
For the last assertion,
$L^{-1}\int_u^{u+L}\|R\|^2\leq L^{-1}\|R\|_{L^2}^2$ uniformly over
allowed origins, and the same Cauchy--Schwarz argument applies.
\end{proof}

Lemmas~\ref{lem:pf-bins} and \ref{lem:pf-transfer} apply to the vector
in \eqref{eq:pf-deletion}.  Its Fourier coefficients are exactly
$z(\gamma)$ for $\gamma\notin J\cup\{0\}$.  The decay of the bin
masses proves square summability of these coefficients.  Lemma
\ref{lem:pf-moments} therefore establishes \eqref{eq:pf-moments} and
\eqref{eq:pf-cov}, including their uniformity.  Pairing opposite
frequencies makes $C_J$ real symmetric.  Positive \emph{definiteness}
requires arithmetic input, which we give next.

For a fixed real $u\geq2$, the fixed-argument Landau formula under
Dirichlet GRH is
\begin{equation}
\label{eq:pf-landau}
 \sum_{T<\gamma\leq T+H}m_\chi(\gamma)e^{i\gamma\log u}
 =-\frac{H}{2\pi\sqrt{u}}\Lambda(u)\chi^*(u)
                       +O_{q,u}(\log(2T)),
 \qquad 1\leq H\leq T.
\end{equation}
Here $\Lambda(u)\chi^*(u)$ is zero unless $u$ is an integral prime
power.  This is the fixed-argument specialization of the
Landau--Gonek formula for Dirichlet $L$-functions
\cite[Theorem~1.2]{Banks2023}; it is also the input to the article's
multiscale coercivity proof \cite{ProveItArticle}.  At a fixed
non-prime-power argument the distance to the neighboring prime powers
is positive.  Endpoints change the sum by at most $O_q(\log(2T))$, by
\eqref{eq:pf-localcount}.

\begin{lemma}[Prime probes]
\label{lem:pf-probes}
For each $a\in G$ one can choose a prime $p_a\equiv a\pmod q$ such
that all numbers $p_a e^{h_k-h_j}$ are at least $2$, and are integral
prime powers only when $j=k$.
\end{lemma}

\begin{proof}
For fixed $R>0$, $R\ne1$, only finitely many primes $p$ make $Rp$ an
integral prime power.  If $R$ is irrational there are none.  Otherwise
write $R=b/c$ in lowest terms.  If $c>1$, integrality forces $c\mid p$
and leaves at most one prime.  If $c=1$, then $b\geq2$ and $bp$ can
be a prime power only if $p$ divides $b$.  There are finitely many
off-diagonal ratios $e^{h_k-h_j}\ne1$.  Exclude their exceptional
primes, take $p$ large enough for the lower bound $2$, and use
Dirichlet's theorem in the reduced class $a$.
\end{proof}

\begin{lemma}[An unsmoothed band is coercive]
\label{lem:pf-coercivity}
There are constants $c,K,T_0>0$, depending only on $q$ and the shifts,
such that, for $T\geq T_0$ and $K\log(2T)\leq H\leq T$,
\begin{equation}
\label{eq:pf-coercivity}
 2\sum_{\gamma\in\Gamma_q\cap(T,T+H]}
 \frac{|A_w(\gamma)|^2}{\tfrac14+\gamma^2}
       \geq\frac{cH}{T^2\log(2T)}\|w\|_2^2,
 \quad
 A_w(\gamma)=\sum_{j,a}w_{j,a}e^{i\gamma h_j}v_a(\gamma).
\end{equation}
This holds for complex coefficient arrays $w$ as well as real ones.
\end{lemma}

\begin{proof}
Let $B=\Gamma_q\cap(T,T+H]$, and use the primes from
Lemma~\ref{lem:pf-probes}.  Expanding \eqref{eq:pf-v} and applying
\eqref{eq:pf-landau} at the finitely many fixed arguments
$p_a e^{h_k-h_j}$ gives
\begin{align}
\label{eq:pf-recover}
 \sum_{\gamma\in B}A_w(\gamma)
              e^{i\gamma(\log p_a-h_j)}
 &=-\frac{H\log p_a}{2\pi\sqrt{p_a}}w_{j,a}
                    +O_{q,h}\bigl(\|w\|_2\log(2T)\bigr).
\end{align}
All off-diagonal main terms vanish.  On the diagonal,
$\chi^*(p_a)=\chi(a)$ because $p_a$ is prime to $q$, and character
orthogonality isolates $w_{j,a}$.

Choose a coordinate with $|w_{j,a}|\geq\|w\|_2/\sqrt D$.
The positive number
$\min_a\log p_a/(2\pi\sqrt{p_a})$ is fixed.  By taking $K$ large
enough, the absolute value on the left of \eqref{eq:pf-recover} is
at least $c_1H\|w\|_2$.  The local zero count, summed over the band,
gives $|B|\leq c_2H\log(2T)$.  Cauchy--Schwarz over its
\emph{distinct} ordinates therefore proves
\[
 \sum_{\gamma\in B}|A_w(\gamma)|^2
                 \geq c_3\frac{H}{\log(2T)}\|w\|_2^2.
\]
Since $\tfrac14+\gamma^2\leq5T^2$ in this band for $T\geq1$,
division by the denominator proves \eqref{eq:pf-coercivity}.
\end{proof}

For real $w$, pairing $\gamma$ and $-\gamma$ in
\eqref{eq:pf-cov} shows that $w^{\mathsf T}C_Jw$ is at least the
left side of \eqref{eq:pf-coercivity} for every band above the deleted
frequencies.  This proves \eqref{eq:pf-band} and positive
definiteness, including after arbitrary fixed finite deletion.  The
argument separates scales by prime powers; it does not separate
individual nearby zeros.

\subsection{Two-sided occupation and failure of exact linear prediction}
\label{subsec:pf-occupation}

We complete the proof of Theorem~\ref{thm:pf-main} by a moment
argument which does not require pointwise boundedness of the
unsmoothed errors.

\begin{lemma}[Occupation from a fourth-moment bound]
\label{lem:pf-occupation}
Suppose a real vector-valued function $Y$ has uniform long-window mean
zero and a uniform long-window second moment $C$ with
$\lambda=\lambda_{\min}(C)>0$.  Suppose also that, for every real
unit $w$ and every sufficiently long allowed interval $I$,
$|I|^{-1}\int_I|w\cdot Y|^4\leq M^4$, where $M\geq1$ is independent
of $w$ and $I$.  Then positive and negative directional occupation
hold uniformly, with permissible constants
\begin{equation}
\label{eq:pf-constants}
 a=\frac{(\lambda/2)^{3/2}}{M^2},
 \qquad \varepsilon=\frac a8,
 \qquad \eta=\left(\frac{a}{8M}\right)^2.
\end{equation}
\end{lemma}

\begin{proof}
Let $\mathbb E$ denote average over the interval and set $f=w\cdot Y$.
Uniformity of the first two moments allows us to make
$\mathbb Ef^2\geq\lambda/2$ and $|\mathbb Ef|\leq a/2$
simultaneously for all unit $w$.  Interpolation between $L^1$ and
$L^4$, equivalently H\"older's inequality, gives
\[
 \mathbb E|f|\geq
       \frac{(\mathbb Ef^2)^{3/2}}{(\mathbb Ef^4)^{1/2}}
       \geq a.
\]
Writing $f_+=\max(f,0)$ and $f_-=\max(-f,0)$, we obtain
$\mathbb Ef_+\geq a/4$ and $\mathbb Ef_-\geq a/4$.
With $\varepsilon=a/8$ this implies
\[
 \frac a8\leq\mathbb E\bigl[f_+\mathbf1_{\{f\geq\varepsilon\}}\bigr]
       \leq(\mathbb Ef^2)^{1/2}
                       \mathbb P(f\geq\varepsilon)^{1/2}
       \leq M\mathbb P(f\geq\varepsilon)^{1/2}.
\]
Here $\mathbb P$ is relative Lebesgue measure on the interval.
The same calculation for $-f$ proves both assertions with the
specified $\eta$.
\end{proof}

The Fourier vector obtained from \eqref{eq:pf-F} by shifts and finite
deletion has bounded $S^4$ norm.  The remainder from
Lemma~\ref{lem:pf-transfer} has bounded local fourth moments too.
Covering a long interval by unit intervals therefore supplies a
common $M$ for $Y=\mathcal X_J$.  The covariance and mean were
established uniformly above, so Lemma~\ref{lem:pf-occupation} proves
\eqref{eq:pf-positive}--\eqref{eq:pf-negative}.  If the closed
convex hull of the essential values failed to contain the
$\varepsilon$-ball, separation would give a unit direction with
essential supremum less than $\varepsilon$, contradicting
\eqref{eq:pf-positive}.  This completes the main theorem.

\begin{corollary}[No exact finite linear explanation]
\label{cor:pf-prediction}
For a nonzero real array $w=(w_{j,a})$, an arbitrary real finite
trigonometric polynomial $P$, and an arbitrary real constant $b$,
the long-window mean square of
\[
 \sum_{j,a}w_{j,a}X_a(t+h_j)-P(t)-b
\]
has a strictly positive limit, uniformly over the allowed origins.
In particular, no nonzero contrast at one fixed scale can be
predicted exactly in this mean-square sense by real linear
combinations of the other finitely many scales and a finite
trigonometric correction.
\end{corollary}

\begin{proof}
Let $J$ contain all nonzero frequencies of $P$ which belong to
$\Gamma_q$, together with their negatives.  Frequencies outside
this finite set retain precisely the coefficients of
$w\cdot\mathcal X_J$.  The uniform Fourier mean-square calculation
of Lemma~\ref{lem:pf-moments}, including any extra frequencies of
$P$, shows that the residual mean square is at least
$w^{\mathsf T}C_Jw>0$.  For a predictor, collect all coefficients in
one array; its block at the predicted scale is the prescribed
nonzero contrast, so the full array cannot vanish.
\end{proof}

This corollary concerns every \emph{fixed} finite correction.  It
does not give a positive lower bound uniform over all corrections:
successively larger Fourier truncations do approximate the error
in mean square.  Likewise, the occupation theorem describes
individual linear contrasts of centered counts.  Positive
covariance and two-sided directional occupation alone do not force
the limiting values to fill an open set, and therefore do not
imply inclusivity of every ordinary prime race.
